% The swept-face advection. Panels a, b: fig_sweptflux_a.tex, fig_sweptflux_b.tex, written by
% make_fig_sweptflux.py (run on Lichtenberg). Panel c: a two-dimensional cut, cells of unit size;
% the face f = {x = 1, 0 <= y <= 1}, displacement (0.5, 0.4), pre-image {x = 0.5, -0.4 <= y <= 0.6};
% the donor D = [0,1]x[0,1] with its plane n = (0.371, 0.928), n.x = 0.603 (through (0.5, 0.45) and
% (1, 0.25)); the lateral neighbour N = [0,1]x[-1,0] full. Closure of D's liquid piece:
% 0.25 (1,0) + 0.45 (-1,0) + 0.5 (0,-1) + 0.5385 (0.371, 0.928) = 0; of N's piece:
% 0.4 (-1,0) + (0.4,-0.5) + 0.5 (0,1) = 0. Liquid: 0.175 + 0.1 of the swept 0.5; alpha_f = 0.85.
\begin{figure}[tb]
\centering
\begin{tikzpicture}[>=Stealth, font=\footnotesize, line join=round]
  \colorlet{facecol}{orange!80!black}
  \colorlet{precol}{teal!60!black}
  \colorlet{wetcol}{yellow!70!orange}
  \colorlet{capcol}{blue!55}
  \colorlet{ifacecol}{blue!65!black}
  \colorlet{fluidcol}{blue!18}
  \colorlet{gaincol}{blue!45}
  \colorlet{bdrycol}{black!50}
  % ------------------------------------------------ (a) the closed triangulation
  \begin{scope}[scale=2.45]
    % written by make_fig_sweptflux.py -- do not edit. d = (0.55, 0.3, 0.22), n = (0.28222, 0.18814, 0.94072), x0 = (0.75, 0.35, 0.5)
% signed volume 0.550000000000 (d.S_f = 0.550000000000), closure |sum a_T| = 0.0e+00, V^- = 0.339625000000, midpoint quadrature (160^3) 0.339625,
% cap area by closure 0.584658019700, by its polygon 0.584658019700
\draw[facecol!60, line width=0.3pt] (1.2000,0.8500) -- (0.7000,0.3500) -- (1.7000,0.3500) -- cycle;
\draw[precol!70, line width=0.3pt] (0.5150,0.4375) -- (1.0150,-0.0625) -- (0.0150,-0.0625) -- cycle;
\draw[black!35, line width=0.3pt] (0.7000,0.3500) -- (0.0150,-0.0625) -- (1.0150,-0.0625) -- cycle;
\draw[black!35, line width=0.3pt] (0.7000,0.3500) -- (1.0150,-0.0625) -- (1.7000,0.3500) -- cycle;
\draw[facecol!60, line width=0.3pt] (1.2000,0.8500) -- (1.7000,0.3500) -- (1.7000,1.3500) -- cycle;
\draw[precol!70, line width=0.3pt] (0.5150,0.4375) -- (1.0150,0.9375) -- (1.0150,-0.0625) -- cycle;
\draw[black!35, line width=0.3pt] (1.7000,0.3500) -- (1.0150,-0.0625) -- (1.0150,0.9375) -- cycle;
\draw[black!35, line width=0.3pt] (1.7000,0.3500) -- (1.0150,0.9375) -- (1.7000,1.3500) -- cycle;
\draw[facecol!60, line width=0.3pt] (1.2000,0.8500) -- (1.7000,1.3500) -- (0.7000,1.3500) -- cycle;
\draw[precol!70, line width=0.3pt] (0.5150,0.4375) -- (0.0150,0.9375) -- (1.0150,0.9375) -- cycle;
\draw[black!35, line width=0.3pt] (1.7000,1.3500) -- (1.0150,0.9375) -- (0.0150,0.9375) -- cycle;
\draw[black!35, line width=0.3pt] (1.7000,1.3500) -- (0.0150,0.9375) -- (0.7000,1.3500) -- cycle;
\draw[facecol!60, line width=0.3pt] (1.2000,0.8500) -- (0.7000,1.3500) -- (0.7000,0.3500) -- cycle;
\draw[precol!70, line width=0.3pt] (0.5150,0.4375) -- (0.0150,-0.0625) -- (0.0150,0.9375) -- cycle;
\draw[black!35, line width=0.3pt] (0.7000,1.3500) -- (0.0150,0.9375) -- (0.0150,-0.0625) -- cycle;
\draw[black!35, line width=0.3pt] (0.7000,1.3500) -- (0.0150,-0.0625) -- (0.7000,0.3500) -- cycle;
\draw[facecol, line width=1.4pt] (0.7000,0.3500) -- (1.7000,0.3500) -- (1.7000,1.3500) -- (0.7000,1.3500) -- cycle;
\draw[precol, line width=1.2pt, dashed] (0.0150,-0.0625) -- (1.0150,-0.0625) -- (1.0150,0.9375) -- (0.0150,0.9375) -- cycle;
\draw[->, black!65, line width=0.6pt] (0.7000,0.3500) -- (0.0561,-0.0378);
\draw[->, black!65, line width=0.6pt] (1.7000,0.3500) -- (1.0561,-0.0378);
\draw[->, black!65, line width=0.6pt] (1.7000,1.3500) -- (1.0561,0.9623);
\draw[->, black!65, line width=0.6pt] (0.7000,1.3500) -- (0.0561,0.9623);
\draw[->, black!65, line width=0.6pt] (1.2000,0.8500) -- (0.5561,0.4622);
\fill[facecol] (1.2000,0.8500) circle[radius=0.018];
\fill[precol] (0.5150,0.4375) circle[radius=0.018];
\coordinate (fC) at (1.2000,0.8500);
\coordinate (pC) at (0.5150,0.4375);
\coordinate (fS) at (1.5150,1.0075);
\coordinate (fP) at (1.7000,1.3500);
\coordinate (pP) at (0.0150,-0.0625);

    \draw[->, facecol, line width=0.9pt] (fC) -- ($(fC)!2.1!(fS)$) node[right=0pt] {$\mathbf{S}_f$};
    \node[facecol, anchor=south west] at (fP) {$f$};
    \node[precol, anchor=north east] at (pP) {$f'$};
    \node[anchor=north] at (0.9,-0.18) {(a)};
  \end{scope}
  % ------------------------------------------------ (b) the cut by the donor's plane
  \begin{scope}[xshift=5.0cm, scale=2.45]
    % written by make_fig_sweptflux.py -- do not edit. d = (0.55, 0.3, 0.22), n = (0.28222, 0.18814, 0.94072), x0 = (0.75, 0.35, 0.5)
% signed volume 0.550000000000 (d.S_f = 0.550000000000), closure |sum a_T| = 0.0e+00, V^- = 0.339625000000, midpoint quadrature (160^3) 0.339625,
% cap area by closure 0.584658019700, by its polygon 0.584658019700
\draw[black!30, line width=0.3pt] (1.2000,0.8500) -- (0.7000,0.3500) -- (1.7000,0.3500) -- cycle;
\draw[black!30, line width=0.3pt] (0.5150,0.4375) -- (1.0150,-0.0625) -- (0.0150,-0.0625) -- cycle;
\draw[black!30, line width=0.3pt] (0.7000,0.3500) -- (0.0150,-0.0625) -- (1.0150,-0.0625) -- cycle;
\draw[black!30, line width=0.3pt] (0.7000,0.3500) -- (1.0150,-0.0625) -- (1.7000,0.3500) -- cycle;
\draw[black!30, line width=0.3pt] (1.2000,0.8500) -- (1.7000,0.3500) -- (1.7000,1.3500) -- cycle;
\draw[black!30, line width=0.3pt] (0.5150,0.4375) -- (1.0150,0.9375) -- (1.0150,-0.0625) -- cycle;
\draw[black!30, line width=0.3pt] (1.7000,0.3500) -- (1.0150,-0.0625) -- (1.0150,0.9375) -- cycle;
\draw[black!30, line width=0.3pt] (1.7000,0.3500) -- (1.0150,0.9375) -- (1.7000,1.3500) -- cycle;
\draw[black!30, line width=0.3pt] (1.2000,0.8500) -- (1.7000,1.3500) -- (0.7000,1.3500) -- cycle;
\draw[black!30, line width=0.3pt] (0.5150,0.4375) -- (0.0150,0.9375) -- (1.0150,0.9375) -- cycle;
\draw[black!30, line width=0.3pt] (1.7000,1.3500) -- (1.0150,0.9375) -- (0.0150,0.9375) -- cycle;
\draw[black!30, line width=0.3pt] (1.7000,1.3500) -- (0.0150,0.9375) -- (0.7000,1.3500) -- cycle;
\draw[black!30, line width=0.3pt] (1.2000,0.8500) -- (0.7000,1.3500) -- (0.7000,0.3500) -- cycle;
\draw[black!30, line width=0.3pt] (0.5150,0.4375) -- (0.0150,-0.0625) -- (0.0150,0.9375) -- cycle;
\draw[black!30, line width=0.3pt] (0.7000,1.3500) -- (0.0150,0.9375) -- (0.0150,-0.0625) -- cycle;
\draw[black!30, line width=0.3pt] (0.7000,1.3500) -- (0.0150,-0.0625) -- (0.7000,0.3500) -- cycle;
\fill[wetcol, opacity=0.45] (1.1125,0.7625) -- (0.7000,0.3500) -- (1.7000,0.3500) -- (1.3313,0.7188) -- cycle;
\fill[wetcol!70!orange, opacity=0.45] (0.5150,0.4375) -- (1.0150,-0.0625) -- (0.0150,-0.0625) -- cycle;
\fill[wetcol, opacity=0.45] (0.7000,0.3500) -- (0.0150,-0.0625) -- (1.0150,-0.0625) -- cycle;
\fill[wetcol, opacity=0.45] (0.7000,0.3500) -- (1.0150,-0.0625) -- (1.7000,0.3500) -- cycle;
\fill[wetcol, opacity=0.45] (1.3312,0.7188) -- (1.7000,0.3500) -- (1.7000,0.6450) -- cycle;
\fill[wetcol!70!orange, opacity=0.45] (0.5150,0.4375) -- (0.7983,0.7208) -- (1.0150,0.6775) -- (1.0150,-0.0625) -- cycle;
\fill[wetcol, opacity=0.45] (1.7000,0.3500) -- (1.0150,-0.0625) -- (1.0150,0.6775) -- (1.3359,0.6623) -- cycle;
\fill[wetcol, opacity=0.45] (1.7000,0.3500) -- (1.3359,0.6623) -- (1.7000,0.6450) -- cycle;
\fill[wetcol!70!orange, opacity=0.45] (0.5150,0.4375) -- (0.0900,0.8625) -- (0.7983,0.7208) -- cycle;
\fill[wetcol, opacity=0.45] (0.7000,0.8450) -- (0.7000,0.3500) -- (1.1125,0.7625) -- cycle;
\fill[wetcol!70!orange, opacity=0.45] (0.5150,0.4375) -- (0.0150,-0.0625) -- (0.0150,0.8775) -- (0.0900,0.8625) -- cycle;
\fill[wetcol, opacity=0.45] (0.0150,0.8775) -- (0.0150,-0.0625) -- (0.4606,0.8564) -- cycle;
\fill[wetcol, opacity=0.45] (0.4606,0.8564) -- (0.0150,-0.0625) -- (0.7000,0.3500) -- (0.7000,0.8450) -- cycle;
\fill[capcol, opacity=0.55] (0.0900,0.8625) -- (0.0150,0.8775) -- (0.4606,0.8564) -- (0.7000,0.8450) -- (1.1125,0.7625) -- (1.3313,0.7188) -- (1.7000,0.6450) -- (1.3359,0.6623) -- (1.0150,0.6775) -- (0.7983,0.7208) -- cycle;
\draw[ifacecol, line width=1.0pt] (0.0900,0.8625) -- (0.0150,0.8775) -- (0.4606,0.8564) -- (0.7000,0.8450) -- (1.1125,0.7625) -- (1.3313,0.7188) -- (1.7000,0.6450) -- (1.3359,0.6623) -- (1.0150,0.6775) -- (0.7983,0.7208) -- cycle;
\draw[facecol, line width=1.4pt] (0.7000,0.3500) -- (1.7000,0.3500) -- (1.7000,1.3500) -- (0.7000,1.3500) -- cycle;
\draw[precol, line width=1.2pt, dashed] (0.0150,-0.0625) -- (1.0150,-0.0625) -- (1.0150,0.9375) -- (0.0150,0.9375) -- cycle;
\fill[black, opacity=0.10] (0.5150,0.4375) -- (1.0150,-0.0625) -- (0.0150,-0.0625) -- cycle;
\draw[black!70, dashed, line width=0.5pt] (0.8750,0.7625) -- (0.5150,0.4375);
\draw[black!70, dashed, line width=0.5pt] (0.8750,0.7625) -- (1.0150,-0.0625);
\draw[black!70, dashed, line width=0.5pt] (0.8750,0.7625) -- (0.0150,-0.0625);
\fill[black] (0.8750,0.7625) circle[radius=0.02];
\coordinate (xz) at (0.8750,0.7625);
\coordinate (capC) at (0.8559,0.7628);
\coordinate (capN) at (0.9909,1.1266);
\coordinate (tC) at (0.5150,0.1042);

    \draw[->, red!65!black, line width=0.9pt] (capC) -- (capN) node[above=0pt] {$\mathbf{n}$};
    \node[anchor=north west, inner sep=1pt] at (xz) {$\mathbf{x}_0$};
    \node[ifacecol, anchor=south] at (0.25,1.15) {$\Gamma$};
    \draw[ifacecol, line width=0.3pt] (0.27,1.15) -- (0.55,0.98);
    \node[anchor=north] at (0.9,-0.18) {(b)};
  \end{scope}
  % ------------------------------------------------ (c) the pieces, two-dimensional cut
  \begin{scope}[xshift=10.6cm, yshift=1.75cm, scale=1.75]
    % liquid: D below its plane, R below its own plane, the lower row full
    \fill[fluidcol] (0,-1) rectangle (2,0);
    \fill[fluidcol] (0,0) -- (1,0) -- (1,0.25) -- (0,0.65) -- cycle;
    \fill[fluidcol] (1,0) -- (2,0) -- (2,0.03) -- (1,0.22) -- cycle;
    \draw[bdrycol, line width=0.5pt] (0,-1) grid[step=1] (2,1);
    \draw[ifacecol, line width=1.4pt] (0,0.65) -- (1,0.25);
    \draw[ifacecol, line width=1.4pt] (1,0.22) -- (2,0.03);
    % the swept volume and its two pieces
    \fill[black!10] (1,0) -- (1,1) -- (0.5,0.6) -- (0.5,0) -- cycle;
    \fill[gaincol] (1,0) -- (1,0.25) -- (0.5,0.45) -- (0.5,0) -- cycle;
    \fill[gaincol] (1,0) -- (0.5,0) -- (0.5,-0.4) -- cycle;
    \draw[black!60, line width=0.5pt] (1,0) -- (1,1) -- (0.5,0.6) -- (0.5,-0.4) -- cycle;
    \draw[ifacecol, line width=1.4pt] (0.5,0.45) -- (1,0.25);
    % the wetted parts of the boundary of the swept volume and the face
    \draw[precol, line width=1.2pt, dashed] (0.5,-0.4) -- (0.5,0.6);
    \draw[facecol, line width=2.2pt] (1,0) -- (1,1);
    \draw[wetted] (0.5,-0.4) -- (0.5,0.45);
    \draw[wetted] (1,0) -- (1,0.25);
    \draw[wetted] (0.5,-0.4) -- (1,0);
    % the caps and their normals: the cell plane y = 0 (both pieces) and D's plane
    \draw[->, black!75, line width=0.7pt] (0.8,0.0) -- (0.8,-0.22) node[right=-1pt] {$A_1\mathbf{n}_1$};
    \draw[->, black!75, line width=0.7pt] (0.62,0.0) -- (0.62,0.22) node[left=-1pt, pos=0.9] {$A_3\mathbf{n}_3$};
    \draw[->, red!65!black, line width=0.7pt] (0.75,0.35) -- ++(0.0928,0.232) node[above left=-2pt] {$A_2\mathbf{n}_2$};
    \node at (0.15,0.88) {$D$};
    \node at (0.15,-0.85) {$N$};
    \node at (1.82,0.85) {$R$};
    \node[facecol, anchor=west] at (1.02,0.88) {$f$};
    \node[precol, anchor=east] at (0.48,-0.28) {$f'$};
    \node[anchor=north] at (1,-1.08) {(c)};
  \end{scope}
\end{tikzpicture}
\caption{The swept-face advection. (a) A face $f$ (orange) with its area vector $\mathbf{S}_f$, traced backward to its pre-image $f'$ (dashed) along the trajectories of its vertices and centre (arrows); the swept volume is closed by the face fan, the reversed pre-image fan and two triangles per lateral quadrilateral, $4F = 16$ triangles for a quadrilateral face. (b) The same volume cut by a plane with normal $\mathbf{n}$: the wetted parts of the triangles (gold), the cap $\Gamma$ (blue), and the cone from the point $\mathbf{x}_0$ of the plane to one wetted triangle (dashed), one term of \cref{eq:swept-volume}; drawn from the geometry computed by \texttt{make\_fig\_sweptflux.py}, whose closed form reproduces a quadrature of the liquid indicator. (c) A two-dimensional cut of the pieces: the swept volume of $f$ reaches from its donor $D$, an interface cell clipped by its own plane, into the full neighbour $N$, wet in total; each liquid piece is closed by the wetted parts of the swept volume's boundary (gold) and by its caps on the cell plane and on $D$'s plane, whose areas follow from the closure \cref{eq:swept-closure}. The wetted part of $f'$ gives $\alpha_f^{n+1} = 0.85$.}
\label{fig:sweptflux}
\end{figure}
